% ==================== 2.4 本章小结 ====================
\section{本章小结}
\label{sec:ch2_summary}

本章针对经典数字孪生五维模型在月面巡视器自主行走场景下的局限性，提出了增强的五维数字孪生模型，并构建了总体架构。主要工作包括：

（1）分析了经典数字孪生五维模型的核心内涵，指出了其在物理实体维度缺乏环境融合、虚拟实体维度缺乏实时演化机制、服务维度缺乏自主决策支持、孪生数据维度缺乏多源异构融合、连接维度缺乏时延补偿机制等五方面的局限性。

（2）提出了面向月面巡视器自主行走的增强五维模型，从物理实体、虚拟实体、服务、孪生数据、连接五个维度进行了适应性增强。增强的物理实体将环境因素纳入考量，形成"巡视器-环境"一体化系统；增强的虚拟实体具备自主演化能力；增强的服务提供自主决策支持；增强的孪生数据实现多源异构融合；增强的连接引入时延补偿机制。

（3）构建了月面巡视器自主行走数字孪生的总体架构，采用分层设计思想，自下而上分为数据层、模型层、服务层和应用层，明确了各层的功能和接口关系，为后续章节的研究提供了统一框架。

本章提出的增强五维模型和总体架构，为月面巡视器自主行走数字孪生系统的构建奠定了理论基础和框架支撑。
